\section{Expressive Power of Knowledge Graph Embeddings [10 points]}
\noindent\textit{Coverage: This problem is covered in Lecture 10.}

TransE is a common method for learning representations of entities and relations in a knowledge graph. Given a triplet $(h, \ell, t)$, where entities embedded as $h$ and $t$ are related by a relation embedded as $\ell$, TransE trains entity and relation embeddings to make $h+\ell$ close to $t$. There are some common patterns that relations form:
\begin{itemize}
    \item Symmetry: A is married to B, and B is married to A.
    \item Inverse: A is teacher of B, and B is student of A. Note that teacher and student are 2 different relations and have their own embeddings.
    \item Composition: $\mathrm{A}$ is son of $\mathrm{B} ; \mathrm{C}$ is sister of $\mathrm{B}$, then $\mathrm{C}$ is aunt of $\mathrm{A}$. Again note that son, sister, and aunt are 3 different relations and have their own embeddings.
\end{itemize}
\subsection{TransE Modeling [3 points]}
For each of the above relational patterns, can TransE model it perfectly, such that $h+\ell=t$ for all relations? Explain why or why not. Note that here $\mathbf{0}$ embeddings for relation are undesirable since that means two entities related by that relation are identical and not distinguishable.

\Solution{
    \\
    \begin{enumerate}
        \item TransE can't model relation of symmetry, since $$\mathbf{h}+\boldsymbol{\ell}=\mathbf{t}$$ $$\mathbf{t}+\boldsymbol{\ell}=\mathbf{h}+2\boldsymbol{\ell} \neq \mathbf{h}$$
        \item TransE can model relation of inverse, for example set $\boldsymbol{\ell}_1=-\boldsymbol{\ell}_2$, when $$\mathbf{h}+\boldsymbol{\ell}_1=\mathbf{t}$$ $$\mathbf{t}+\boldsymbol{\ell}_2=\mathbf{h}$$
        \item TransE can model relation of composition, set $\boldsymbol{\ell}_3=\boldsymbol{\ell}_1+\boldsymbol{\ell}_2$, when $$\mathbf{a}+\boldsymbol{\ell}_1=\mathbf{b}$$ $$\mathbf{b}+\boldsymbol{\ell}_2=\mathbf{c}$$, we have $$\mathbf{a}+\boldsymbol{\ell}_3=\mathbf{c}$$
    \end{enumerate}
}

\subsection{RotatE Modeling [3 points]}
Consider a new model, RotatE. Instead of training embeddings such that $h+\ell \approx t$, we train embeddings such that $h \circ \ell \approx t$. Here $\circ$ means rotation. You can think of $h$ as a vector of dimension $2 d$, representing $d$ $2 \mathrm{D}$ points. $\ell$ is a $d$-dimensional vector specifying rotation angles. When applying $\circ$, For all $i \in 0 \ldots d-1, \left(h_{2 i}, h_{2 i+1}\right)$ is rotated clockwise by $l_i$. Similar to TransE, the entity embeddings are also normalized to L2 norm 1. Can RotatE model the above 3 relation patterns perfectly such that $h \circ \ell = t$ for all relations? Why or why not?

\Solution{
    \\
    \begin{enumerate}
        \item RotatE can model relation of symmetry, set $l_i=2k\pi,\enspace k \in \mathbb{Z}$, so $$\mathbf{h}\circ\boldsymbol{\ell}=\mathbf{t}$$ $$\mathbf{t}\circ\boldsymbol{\ell}=\mathbf{h}$$
        \item RotatE can model relation of inverse, set $l^{(2)}_i=-l^{(1)}_i+ 2k\pi,\enspace k \in \mathbb{Z}$, so $$\mathbf{h}\circ\boldsymbol{\ell}_1=\mathbf{t}$$ $$\mathbf{t}\circ\boldsymbol{\ell}_2=\mathbf{h}$$
        \item RotatE can model relation of inverse, set $l^{(3)}_i=l^{(1)}_i+l^{(2)}_i+ 2k\pi,\enspace k \in \mathbb{Z}$, when $$\mathbf{a}\circ\boldsymbol{\ell}_1=\mathbf{b}$$ $$\mathbf{b}\circ\boldsymbol{\ell}_2=\mathbf{c}$$, so $$\mathbf{a}\circ\boldsymbol{\ell}_3=\mathbf{c}$$
    \end{enumerate}
}

\subsection{Failure Cases [4 points]}
Give an example of a graph that RotatE cannot model. Can TransE model this graph? Assume that relation embeddings cannot be $\mathbf{0}$ in either model.

\Solution{}